\documentclass[reprint,amsmath,amssymb,aps,prl,floatfix,onecolumn]{revtex4-2}
\usepackage[brazilian]{babel}
\usepackage[T1]{fontenc}
\usepackage{graphicx}
\usepackage{bm}
\usepackage{color}
\usepackage{float}
\usepackage{hyperref}
\usepackage{longtable}

% ============================================================================
%  NOTAS DE EDICAO (apagaveis de uma vez so)
%  --------------------------------------------------------------------------
%  As notas em AMARELO (\nota{...}) e os placeholders de print (\missingfigure)
%  servem so de guia para voce e sua dupla. Para gerar a VERSAO FINAL limpa
%  (sem nenhuma nota nem caixa de print), basta trocar a linha abaixo por:
%        \usepackage[disable]{todonotes}
%  Com [disable], TODAS as notas e placeholders somem do PDF automaticamente.
% ============================================================================
\usepackage{todonotes}
\newcommand{\nota}[1]{\todo[inline,color=yellow!55,bordercolor=orange,size=\small]{\textbf{NOTA:} #1}}

\begin{document}

% A capa e opcional; e o seu arquivo capa.tex (suba-o junto no Overleaf).
% \input{capa.tex}

\title{Visualizador Georreferenciado de Segurança Pública do Estado de São Paulo}
\author{Giovanni Monteiro Franzini}
\email{monteirogmf@ita.br}
\affiliation{Instituto Tecnológico de Aeronáutica, São José dos Campos, Brasil}
\date{\today}

\begin{abstract}
\textit{Este trabalho apresenta um visualizador web interativo de dados georreferenciados
para a análise estatística e espaço-temporal das ocorrências criminais no Estado de São Paulo,
no período de janeiro de 2022 a abril de 2026, a partir dos dados abertos oficiais da Secretaria
de Segurança Pública do Estado de São Paulo (SSP-SP). A aplicação, construída com OpenLayers,
React e um banco de dados espacial PostgreSQL/PostGIS, processa cerca de 5,17 milhões de registros
e os disponibiliza por meio de mapa de calor, agrupamento por tipo de crime, navegação temporal e
ferramentas de análise de vizinhança (buffer), mantendo uma abordagem estritamente estatística e
neutra.}
\end{abstract}

\maketitle

\nota{Confiram o nome da dupla (incluir o(a) colega no \texttt{\textbackslash author}) e o e-mail. O nome do arquivo final, conforme a professora, deve seguir o padrão \texttt{NomeApp\_Relatorio.pdf} -- escolham um nome curto para o app.}

% ============================================================================
\section{Declaração de Uso de Inteligência Artificial}
% ============================================================================
Declaro que foi feito uso de Inteligência Artificial (assistente Claude) na realização deste
trabalho, de forma detalhada nas seguintes frentes: (i) geração e revisão do código dos
\textit{scripts} de extração, limpeza e carga dos dados (ETL), da API \textit{backend} e do
\textit{frontend} do visualizador; (ii) apoio às decisões de arquitetura de \textit{software}
(escolha do padrão de três camadas, do banco espacial e das estratégias de renderização);
(iii) automação do processo de publicação na nuvem (\textit{deploy}); (iv) diagnóstico e correção
de erros; e (v) estruturação e aprimoramento textual desta documentação. Todo o conteúdo gerado foi
\textbf{conferido, testado e validado} pelos autores, que mantiveram o controle sobre o escopo, as
decisões e a verificação dos resultados.

\nota{Ajustem este parágrafo para refletir exatamente o que vocês fizeram manualmente vs. o que foi assistido por IA, conforme a honestidade de cada um. A professora pede a descrição detalhada do uso.}

% ============================================================================
\section{Visão Geral do Visualizador}
% ============================================================================

\subsection{Área de Aplicação e Motivação}
A área de aplicação escolhida foi a \textbf{Segurança Pública}. A motivação central é transformar
os boletins de ocorrência brutos, disponibilizados em planilhas pela SSP-SP, em uma ferramenta
visual e interativa que permita compreender a \textbf{dinâmica espacial e a evolução temporal} das
manchas criminais no Estado de São Paulo, evidenciando tendências e o possível deslocamento da
criminalidade ao longo dos anos. Em conformidade com as instruções da disciplina, adotou-se uma
abordagem \textbf{estritamente estatística e neutra}, focada no dado e no padrão espacial, sem
qualquer análise política ou social.

\subsection{O que é Visualizado}
A aplicação exibe, sobre um mapa do Estado de São Paulo, a localização das ocorrências criminais.
A visualização é adaptativa: em escala estadual, apresenta um \textbf{mapa de calor} (densidade);
ao aproximar, transiciona para \textbf{pontos individuais agrupados (clusters) coloridos por
natureza criminal}. O usuário pode navegar no tempo (mês a mês, de jan/2022 a abr/2026), filtrar
por município e por tipo de crime, e delimitar áreas de vizinhança. O foco de contingência do
estudo é o município de \textbf{São José dos Campos}.

% ============================================================================
\section{Os Dados}
% ============================================================================

\subsection{Fontes dos Dados}
Os dados utilizados são oriundos de fontes oficiais e abertas:
\begin{itemize}
  \item \textbf{Ocorrências criminais:} planilhas de boletins de ocorrência da Secretaria de
  Segurança Pública do Estado de São Paulo (SSP-SP)~\cite{ssp}, recorte de jan/2022 a abr/2026.
  \item \textbf{Malha municipal:} limites territoriais dos 645 municípios de SP, obtidos via API do
  Instituto Brasileiro de Geografia e Estatística (IBGE)~\cite{ibge}.
  \item \textbf{Centroides de bairro:} coordenadas de bairros obtidas do projeto colaborativo
  OpenStreetMap (via Overpass e Nominatim)~\cite{osm}, para georreferenciar ocorrências sem
  coordenada exata.
\end{itemize}

\subsection{Tratamento dos Dados (ETL)}
As planilhas brutas (cerca de 5,17 milhões de registros, $\sim$1,87 GB) passaram por um processo de
\textit{ETL} (Extração, Transformação e Carga) que tratou anomalias reais encontradas na base:
codificação de caracteres distinta entre os anos, nomes de colunas que mudam ao longo do tempo,
campos com separadores embutidos, coordenadas inválidas e registros sob sigilo legal. A
geolocalização final foi distribuída em três níveis de precisão, resumidos na Tabela~\ref{tab:cobertura}.

\begin{table}[H]
\centering
\caption{Cobertura geográfica final das ocorrências.}
\label{tab:cobertura}
\begin{tabular}{lcc}
\hline
\textbf{Categoria} & \textbf{Ocorrências} & \textbf{\%} \\ \hline
Coordenada exata (original)        & 3.706.674 & 71,7\% \\
Aproximado (centroide de bairro)   & 845.217   & 16,4\% \\
Sem coordenada (oculto no mapa)    & 616.211   & 11,9\% \\ \hline
\textbf{Total}                     & \textbf{5.168.102} & \textbf{100\%} \\ \hline
\end{tabular}
\end{table}
\nota{Se quiserem, posso detalhar mais o ETL (encoding misto, fuzzy, etc.) -- decidam o nível de profundidade adequado ao relatório.}

% ============================================================================
\section{A Interface do Visualizador}
% ============================================================================
A seguir, descrevem-se os principais elementos da interface, acompanhados de capturas de tela.

\missingfigure[figwidth=0.8\linewidth]{PRINT: tela inicial do visualizador (heatmap estadual). Capturar a visão geral com a barra lateral.}

\subsection{Mapa de Calor e Renderização Adaptativa}
Na visão estadual, a densidade de ocorrências é representada por um mapa de calor cuja escala de
intensidade se ajusta à área exibida. Ao aproximar (zoom), o sistema transiciona automaticamente
para pontos individuais agrupados em \textit{clusters}, coloridos segundo a natureza criminal.

\missingfigure[figwidth=0.8\linewidth]{PRINT: visão aproximada de uma cidade mostrando os clusters coloridos por natureza.}

\subsection{Controle Temporal}
Uma barra inferior permite navegar no tempo, com reprodução automática (\textit{play}) e dois modos:
\textbf{Atual} (exibe apenas o mês selecionado) e \textbf{Acumulado} (exibe o intervalo entre um
mês inicial e um final, ambos escolhidos pelo usuário).

\missingfigure[figwidth=0.7\linewidth]{PRINT: a barra do tempo, destacando os dois cursores (início e fim) no modo Acumulado.}

\subsection{Filtros e Ferramenta de Buffer}
A barra lateral oferece filtros por município e por natureza criminal (agrupada por famílias de
cores), além de um botão para incluir as ocorrências de localização aproximada. A ferramenta de
\textbf{buffer} permite expandir a análise para uma área de vizinhança (município + raio em km):
por exemplo, São José dos Campos com 15~km de raio passa a abranger cidades vizinhas como Jacareí.

\missingfigure[figwidth=0.8\linewidth]{PRINT: a ferramenta de buffer aplicada (contorno da área e ocorrências da vizinhança).}

\subsection{Interatividade}
Ao clicar em um ponto, um balão (\textit{popup}) exibe os metadados da ocorrência: mês/ano,
natureza criminal e o nível de precisão geográfica (localização exata ou aproximada).

\missingfigure[figwidth=0.6\linewidth]{PRINT: um popup aberto sobre uma ocorrência.}

% ============================================================================
\section{Arquitetura e Ferramentas}
% ============================================================================
A aplicação adota o padrão \textbf{Web GIS de três camadas}: uma camada de apresentação (frontend),
uma de lógica (API) e uma de dados (banco espacial), conteinerizadas com Docker. A
Tabela~\ref{tab:ferramentas} lista as ferramentas utilizadas, com suas versões, conforme exigido.

\begin{table}[H]
\centering
\caption{Ferramentas utilizadas e respectivas versões.}
\label{tab:ferramentas}
\begin{tabular}{lll}
\hline
\textbf{Camada} & \textbf{Ferramenta} & \textbf{Versão} \\ \hline
Banco de dados   & PostgreSQL          & 16 \\
                 & PostGIS             & 3.4 \\
ETL / Backend    & Node.js             & 22 \\
                 & TypeScript          & 5.7 \\
                 & Fastify             & 5 \\
Frontend         & React               & 18 \\
                 & OpenLayers          & 10 \\
                 & Vite                & 5 \\
Infraestrutura   & Docker / Compose    & --- \\
                 & Caddy               & 2 \\
                 & Nginx               & 1.x \\
Hospedagem       & Google Cloud (Compute Engine, Ubuntu) & 24.04 \\
Entrega contínua & GitHub Actions      & --- \\ \hline
\end{tabular}
\end{table}
\nota{Confiram/atualizem as versões exatas se desejarem precisão máxima (estão nos arquivos \texttt{package.json}). A professora pede explicitamente as versões das ferramentas.}

\subsection{Hospedagem e Publicação}
Toda a aplicação (banco, API e frontend) executa em uma máquina virtual no Google Cloud, com cada
serviço em um \textit{container} Docker. O código é versionado no GitHub e a publicação é
automatizada: a cada envio de código, uma esteira de integração contínua (GitHub Actions) atualiza
o servidor automaticamente. O acesso é feito por endereço seguro (HTTPS).

\missingfigure[figwidth=0.85\linewidth]{PRINT (opcional): diagrama da arquitetura. Pode ser exportado do documento HTML "Como funciona".}

% ============================================================================
\section{Considerações Finais}
% ============================================================================
O projeto demonstrou a viabilidade de transformar uma base oficial massiva e heterogênea em uma
ferramenta de apoio à decisão, visual e interativa, mantendo rigor estatístico. O maior desafio
técnico residiu no tratamento dos dados brutos -- a etapa de ETL --, evidenciando a habitual
dissonância entre os formatos fornecidos pelos órgãos oficiais e a estrutura ideal para um Sistema
de Informação Geográfica. A solução adotada para a geolocalização, que preserva a transparência
estatística sem distorcer a leitura espacial do mapa, ilustra esse tipo de decisão de engenharia.

\nota{Espaço para comentários finais da dupla: aprendizados, dificuldades, possíveis evoluções futuras.}

\begin{thebibliography}{99}
\bibitem{ssp}
SECRETARIA DE SEGURANÇA PÚBLICA DO ESTADO DE SÃO PAULO (SSP-SP). \textbf{Dados Abertos -- Estatísticas}.
Disponível em: \url{https://www.ssp.sp.gov.br/estatistica/consultas}. Acesso em: \today.

\bibitem{ibge}
INSTITUTO BRASILEIRO DE GEOGRAFIA E ESTATÍSTICA (IBGE). \textbf{Malhas Territoriais Municipais}.
Disponível em: \url{https://servicodados.ibge.gov.br/api/docs/malhas}. Acesso em: \today.

\bibitem{osm}
OPENSTREETMAP. \textbf{Dados geográficos colaborativos (Overpass / Nominatim)}.
Disponível em: \url{https://www.openstreetmap.org}. Acesso em: \today.
\end{thebibliography}

\end{document}
